\begin{thebibliography}{999}
\bibitem{planck} M. Planck. \href{https://doi.org/10.1002/andp.19013090310}{Ueber das Gesetz der Energieverteilung im Normalspectrum}. \emph{Annalen der Physik} 309 553-563 (1901).
\bibitem{einstein} A. Einstein. \href{https://doi.org/10.1002/andp.19053220607}{Über einen die Erzeugung und Verwandlung des Lichtes betreffenden heuristischen Gesichtspunkt}. \emph{Annalen der Physik} 322 132-148 (1905).
\bibitem{heisenberg} W. Heisenberg. \href{https://doi.org/10.1007/BF01328377}{Über quantentheoretische Umdeutung kinematischer und mechanischer Beziehungen}. \emph{Zeitschrift für Physik} 33 879-893 (1925).
\bibitem{schrodinger} E. Schrödinger. \href{https://doi.org/10.1103/PhysRev.28.1049}{An Undulatory Theory of the Mechanics of Atoms and Molecules}. \emph{Physical Review} 28 1049-1070 (1926).
\bibitem{born} M. Born. \href{https://doi.org/10.1007/BF01397477}{Zur Quantenmechanik der Stossvorgänge}. \emph{Zeitschrift für Physik} 37 863-867 (1926).
\bibitem{dirac} P. A. M. Dirac. \href{https://doi.org/10.1098/rspa.1927.0012}{The physical interpretation of the quantum dynamics}. \emph{Proceedings of the Royal Society of London. Series A, Containing Papers of a Mathematical and Physical Character} 113 621-641 (1927).
\bibitem{gleason} A. Gleason. \href{https://doi.org/10.1512/iumj.1957.6.56050}{Measures on the Closed Subspaces of a Hilbert Space}. \emph{Indiana University Mathematics Journal} 6 885-893 (1957).
\bibitem{busch} P. Busch. \href{https://doi.org/10.1103/PhysRevLett.91.120403}{Quantum States and Generalized Observables: A Simple Proof of Gleason's Theorem}. \emph{Physical Review Letters} 91 120403 (2003).
\bibitem{qubitgleason} V. Fiorentino, S. Weigert. \href{https://arxiv.org/abs/2511.15607}{Gleason's Theorem for a Qubit as Part of a Composite System}. \emph{arXiv:2511.15607}  (2025).
\bibitem{epr} A. Einstein, B. Podolsky, N. Rosen. \href{https://doi.org/10.1103/PhysRev.47.777}{Can Quantum-Mechanical Description of Physical Reality Be Considered Complete?}. \emph{Physical Review} 47 777-780 (1935).
\bibitem{bell} J. S. Bell. \href{https://doi.org/10.1103/PhysicsPhysiqueFizika.1.195}{On the Einstein Podolsky Rosen paradox}. \emph{Physics Physique Fizika} 1 195-200 (1964).
\bibitem{chsh} J. F. Clauser, M. A. Horne, A. Shimony, R. A. Holt. \href{https://doi.org/10.1103/PhysRevLett.23.880}{Proposed Experiment to Test Local Hidden-Variable Theories}. \emph{Physical Review Letters} 23 880-884 (1969).
\bibitem{ks} S. Kochen, E. Specker. \href{https://doi.org/10.1512/iumj.1968.17.17004}{The Problem of Hidden Variables in Quantum Mechanics}. \emph{Indiana University Mathematics Journal} 17 59-87 (1967).
\bibitem{spekkens} R. W. Spekkens. \href{https://doi.org/10.1103/PhysRevA.71.052108}{Contextuality for preparations, transformations, and unsharp measurements}. \emph{Physical Review A} 71 052108 (2005).
\bibitem{ptm} D. Schmid, R. D. Baldijão, J. H. Selby, A. B. Sainz, R. W. Spekkens. \href{https://doi.org/10.22331/q-2026-10-27-2250}{Noncontextuality inequalities for prepare-transform-measure scenarios}. \emph{Quantum} 10 2250 (2026).
\bibitem{wigner} E. Wigner. \href{https://doi.org/10.1103/PhysRev.40.749}{On the Quantum Correction For Thermodynamic Equilibrium}. \emph{Physical Review} 40 749-759 (1932).
\bibitem{husimi} K. Husimi. \href{https://www.jstage.jst.go.jp/article/ppmsj1919/22/4/22_4_264/_article}{Some Formal Properties of the Density Matrix}. \emph{Proc. Phys.-Math. Soc. Japan 22, 264--314}  (1940).
\bibitem{glauber} R. J. Glauber. \href{https://doi.org/10.1103/PhysRev.131.2766}{Coherent and Incoherent States of the Radiation Field}. \emph{Physical Review} 131 2766-2788 (1963).
\bibitem{sudarshan} E. C. G. Sudarshan. \href{https://doi.org/10.1103/PhysRevLett.10.277}{Equivalence of Semiclassical and Quantum Mechanical Descriptions of Statistical Light Beams}. \emph{Physical Review Letters} 10 277-279 (1963).
\bibitem{wootters} W. K. Wootters. \href{https://doi.org/10.1016/0003-4916(87)90176-X}{A Wigner-function formulation of finite-state quantum mechanics}. \emph{Annals of Physics} 176 1-21 (1987).
\bibitem{frame} C. Ferrie, J. Emerson. \href{https://doi.org/10.1088/1751-8113/41/35/352001}{Frame representations of quantum mechanics and the necessity of negativity in quasi-probability representations}. \emph{Journal of Physics A: Mathematical and Theoretical} 41 352001 (2008).
\bibitem{kd} D. Schmid, R. D. Baldijão, Y. Yīng, R. Wagner, J. H. Selby. \href{https://doi.org/10.1103/PhysRevA.110.052206}{Kirkwood-Dirac representations beyond quantum states and their relation to noncontextuality}. \emph{Physical Review A} 110 052206 (2024).
\bibitem{sic} J. M. Renes, R. Blume-Kohout, A. J. Scott, C. M. Caves. \href{https://doi.org/10.1063/1.1737053}{Symmetric informationally complete quantum measurements}. \emph{Journal of Mathematical Physics} 45 2171-2180 (2004).
\bibitem{bargmann} V. Bargmann. \href{https://doi.org/10.1002/cpa.3160140303}{On a Hilbert space of analytic functions and an associated integral transform part I}. \emph{Communications on Pure and Applied Mathematics} 14 187-214 (1961).
\bibitem{majorana} E. Majorana. \href{https://doi.org/10.1007/BF02960953}{Atomi orientati in campo magnetico variabile}. \emph{Il Nuovo Cimento} 9 43-50 (1932).
\bibitem{holocode} M. W. AlMasri. \href{https://arxiv.org/abs/2609.19162}{Holomorphic Quantum Error Correction Codes}. \emph{arXiv:2609.19162}  (2026).
\bibitem{field} O. Éboli, R. Jackiw, S. Pi. \href{https://doi.org/10.1103/PhysRevD.37.3557}{Quantum fields out of thermal equilibrium}. \emph{Physical Review D} 37 3557-3581 (1988).
\bibitem{feynman} R. P. Feynman. \href{https://doi.org/10.1007/BF02650179}{Simulating physics with computers}. \emph{International Journal of Theoretical Physics} 21 467-488 (1982).
\bibitem{deutsch} D. Deutsch. \href{https://doi.org/10.1098/rspa.1985.0070}{Quantum theory, the Church--Turing principle and the universal quantum computer}. \emph{Proceedings of the Royal Society of London. A. Mathematical and Physical Sciences} 400 97-117 (1985).
\bibitem{bv} E. Bernstein, U. Vazirani. \href{https://doi.org/10.1137/S0097539796300921}{Quantum Complexity Theory}. \emph{SIAM Journal on Computing} 26 1411-1473 (1997).
\bibitem{schumacher} B. Schumacher. \href{https://doi.org/10.1103/PhysRevA.51.2738}{Quantum coding}. \emph{Physical Review A} 51 2738-2747 (1995).
\bibitem{noclone} W. K. Wootters, W. H. Zurek. \href{https://doi.org/10.1038/299802a0}{A single quantum cannot be cloned}. \emph{Nature} 299 802-803 (1982).
\bibitem{dj} D. Deutsch, R. Jozsa. \href{https://doi.org/10.1098/rspa.1992.0167}{Rapid solution of problems by quantum computation}. \emph{Proceedings of the Royal Society of London. Series A: Mathematical and Physical Sciences} 439 553-558 (1992).
\bibitem{shor} P. W. Shor. \href{https://doi.org/10.1137/S0097539795293172}{Polynomial-Time Algorithms for Prime Factorization and Discrete Logarithms on a Quantum Computer}. \emph{SIAM Journal on Computing} 26 1484-1509 (1997).
\bibitem{grover} L. K. Grover. \href{https://doi.org/10.1145/237814.237866}{A fast quantum mechanical algorithm for database search}. \emph{Proceedings of the twenty-eighth annual ACM symposium on Theory of computing - STOC '96} 212-219 (1996).
\bibitem{barenco} A. Barenco, C. H. Bennett, R. Cleve, et al. \href{https://doi.org/10.1103/PhysRevA.52.3457}{Elementary gates for quantum computation}. \emph{Physical Review A} 52 3457-3467 (1995).
\bibitem{boykin} P. O. Boykin, T. Mor, M. Pulver, V. Roychowdhury, F. Vatan. \href{https://doi.org/10.1109/SFFCS.1999.814621}{On universal and fault-tolerant quantum computing: a novel basis and a new constructive proof of universality for Shor's basis}. \emph{40th Annual Symposium on Foundations of Computer Science (Cat. No.99CB37039)} 486-494 (1999).
\bibitem{sk} C. M. Dawson, M. A. Nielsen. \href{https://arxiv.org/abs/quant-ph/0505030}{The Solovay-Kitaev algorithm}. \emph{arXiv:quant-ph/0505030}  (2005).
\bibitem{error} P. W. Shor. \href{https://doi.org/10.1103/PhysRevA.52.R2493}{Scheme for reducing decoherence in quantum computer memory}. \emph{Physical Review A} 52 R2493-R2496 (1995).
\bibitem{teleport} C. H. Bennett, G. Brassard, C. Crépeau, et al. \href{https://doi.org/10.1103/PhysRevLett.70.1895}{Teleporting an unknown quantum state via dual classical and Einstein-Podolsky-Rosen channels}. \emph{Physical Review Letters} 70 1895-1899 (1993).
\bibitem{threshold} Google Quantum AI and Collaborators, R. Acharya, D. A. Abanin, et al. \href{https://doi.org/10.1038/s41586-024-08449-y}{Quantum error correction below the surface code threshold}. \emph{Nature} 638 920-926 (2024).
\bibitem{universalqc} H. Zhou, C. Zhao, M. Cain, et al. \href{https://doi.org/10.1038/s41586-025-09543-5}{Low-overhead transversal fault tolerance for universal quantum computation}. \emph{Nature} 646 303-308 (2025).
\bibitem{category} S. Abramsky, B. Coecke. \href{https://arxiv.org/abs/quant-ph/0402130}{A categorical semantics of quantum protocols}. \emph{arXiv:quant-ph/0402130}  (2004).
\bibitem{jam} A. Jamiolkowski. \href{https://doi.org/10.1016/0034-4877(72)90011-0}{Linear transformations which preserve trace and positive semidefiniteness of operators}. \emph{Reports on Mathematical Physics} 3 275-278 (1972).
\bibitem{choi} M. Choi. \href{https://doi.org/10.1016/0024-3795(75)90075-0}{Completely positive linear maps on complex matrices}. \emph{Linear Algebra and its Applications} 10 285-290 (1975).
\bibitem{stine} W. F. Stinespring. \href{https://doi.org/10.2307/2032342}{Positive Functions on C * -Algebras}. \emph{Proceedings of the American Mathematical Society} 6 211 (1955).
\bibitem{zx} B. Coecke, R. Duncan. \href{https://arxiv.org/abs/0906.4725}{Interacting Quantum Observables: Categorical Algebra and Diagrammatics}. \emph{arXiv:0906.4725}  (2009).
\bibitem{qset} A. Kornell. \href{https://arxiv.org/abs/1804.00581}{Quantum sets}. \emph{arXiv:1804.00581}  (2018).
\bibitem{effectus} K. Cho, B. Jacobs, B. Westerbaan, A. Westerbaan. \href{https://arxiv.org/abs/1512.05813}{An Introduction to Effectus Theory}. \emph{arXiv:1512.05813}  (2015).
\bibitem{qpl} P. Selinger. \href{https://doi.org/10.1017/S0960129504004256}{Towards a quantum programming language}. \emph{Mathematical Structures in Computer Science} 14 527-586 (2004).
\bibitem{lambda} P. Selinger, B. Valiron. \href{https://arxiv.org/abs/cs/0404056}{A lambda calculus for quantum computation with classical control}. \emph{arXiv:cs/0404056}  (2004).
\bibitem{qml} T. Altenkirch, J. Grattage. \href{https://arxiv.org/abs/quant-ph/0409065}{A functional quantum programming language}. \emph{arXiv:quant-ph/0409065}  (2004).
\bibitem{arrows} J. K. Vizzotto, T. Altenkirch, A. Sabry. \href{https://arxiv.org/abs/quant-ph/0501151}{Structuring quantum effects: superoperators as arrows}. \emph{arXiv:quant-ph/0501151}  (2005).
\bibitem{quipper} A. S. Green, P. L. Lumsdaine, N. J. Ross, P. Selinger, B. Valiron. \href{https://arxiv.org/abs/1304.3390}{Quipper: A Scalable Quantum Programming Language}. \emph{arXiv:1304.3390}  (2013).
\bibitem{proto} F. Rios, P. Selinger. \href{https://arxiv.org/abs/1706.02630}{A Categorical Model for a Quantum Circuit Description Language (Extended Abstract)}. \emph{arXiv:1706.02630}  (2017).
\bibitem{qwire} J. Paykin, R. Rand, S. Zdancewic. \href{https://doi.org/10.1145/3009837.3009894}{QWIRE: a core language for quantum circuits}. \emph{Proceedings of the 44th ACM SIGPLAN Symposium on Principles of Programming Languages} 846-858 (2017).
\bibitem{silq} B. Bichsel, M. Baader, T. Gehr, M. Vechev. \href{https://doi.org/10.1145/3385412.3386007}{Silq: a high-level quantum language with safe uncomputation and intuitive semantics}. \emph{Proceedings of the 41st ACM SIGPLAN Conference on Programming Language Design and Implementation} 286-300 (2020).
\bibitem{qunity} F. Voichick, L. Li, R. Rand, M. Hicks. \href{https://arxiv.org/abs/2204.12384}{Qunity: A Unified Language for Quantum and Classical Computing (Extended Version)}. \emph{arXiv:2204.12384}  (2022).
\bibitem{lambdapc} J. Paykin, S. Winnick. \href{https://arxiv.org/abs/2407.16801}{Qudit Quantum Programming with Projective Cliffords}. \emph{arXiv:2407.16801}  (2024).
\bibitem{cpos} A. Kornell, B. Lindenhovius, M. Mislove. \href{https://doi.org/10.1017/S096012952610053X}{Categories of quantum cpos}. \emph{Mathematical Structures in Computer Science} 36 e15 (2026).
\bibitem{super} G. Chiribella, G. M. D'Ariano, P. Perinotti. \href{https://doi.org/10.1209/0295-5075/83/30004}{Transforming quantum operations: Quantum supermaps}. \emph{EPL (Europhysics Letters)} 83 30004 (2008).
\bibitem{combs} G. Chiribella, G. M. D'Ariano, P. Perinotti. \href{https://doi.org/10.1103/PhysRevA.80.022339}{Theoretical framework for quantum networks}. \emph{Physical Review A} 80 022339 (2009).
\bibitem{process} F. A. Pollock, C. Rodríguez-Rosario, T. Frauenheim, M. Paternostro, K. Modi. \href{https://doi.org/10.1103/PhysRevLett.120.040405}{Operational Markov Condition for Quantum Processes}. \emph{Physical Review Letters} 120 040405 (2018).
\bibitem{switch} G. Chiribella, G. M. D'Ariano, P. Perinotti, B. Valiron. \href{https://doi.org/10.1103/PhysRevA.88.022318}{Quantum computations without definite causal structure}. \emph{Physical Review A} 88 022318 (2013).
\bibitem{indef} O. Oreshkov, F. Costa, Č. Brukner. \href{https://doi.org/10.1038/ncomms2076}{Quantum correlations with no causal order}. \emph{Nature Communications} 3 1092 (2012).
\bibitem{projectors} T. Hoffreumon, O. Oreshkov. \href{https://doi.org/10.22331/q-2026-01-21-1978}{Projective characterization of higher-order quantum transformations}. \emph{Quantum} 10 1978 (2026).
\bibitem{jencova} A. Jenčová. \href{https://doi.org/10.22331/q-2026-05-05-2090}{On the structure of higher order quantum maps}. \emph{Quantum} 10 2090 (2026).
\bibitem{jpos} A. Jenčová. \href{https://arxiv.org/abs/2604.09192}{Order structure and signalling in higher order quantum maps}. \emph{arXiv:2604.09192}  (2026).
\bibitem{causaldecomp} T. van der Lugt, R. Lorenz. \href{https://doi.org/10.22331/q-2026-09-30-2218}{Unitary causal decompositions: a characterisation via lattice theory}. \emph{Quantum} 10 2218 (2026).
\bibitem{guide} P. Taranto, S. Milz, M. Murao, M. T. Quintino, K. Modi. \href{https://doi.org/10.1103/92c4-g7qs}{A Guide to Higher-order quantum operations}. \emph{PRX Quantum, accepted June 9}  (2026).
\bibitem{granthi} S. Abramsky, R. Jagadeesan. \href{https://arxiv.org/abs/2608.20443}{Granthi: Higher-Order Quantum Programming via Unitary Wiring}. \emph{arXiv:2608.20443}  (2026).
\bibitem{causaltype} K. Barsse, R. Péchoux, S. Perdrix. \href{https://arxiv.org/abs/2607.09534}{Higher-Order Programs with Indefinite Causal Orders: a Linear Approach to Coherent Control of Quantum Processes}. \emph{arXiv:2607.09534}  (2026).
\bibitem{purecontrol} K. Hirata, T. Tsukada. \href{https://arxiv.org/abs/2607.15926}{Causality in Pure Quantum Computation with Quantum Control}. \emph{arXiv:2607.15926}  (2026).
\bibitem{closedlabs} M. Salzger, V. Vilasini. \href{https://arxiv.org/abs/2605.08351}{Higher-order quantum processes respecting closed labs in a spacetime have quantum controlled causal order}. \emph{arXiv:2605.08351}  (2026).
\bibitem{essential} S. Abramsky, R. Jagadeesan. \href{https://arxiv.org/abs/2606.04080}{Essential Unitarity for Higher-Order Quantum Computation}. \emph{arXiv:2606.04080}  (2026).
\bibitem{quasip} N. D. Ahmed, R. Kunjwal. \href{https://arxiv.org/abs/2610.00579}{Characterizing unitaries via quasi-process functions}. \emph{arXiv:2610.00579}  (2026).
\bibitem{loglambda} K. Tokuo. \href{https://arxiv.org/abs/2610.34266}{A Lambda Calculus Based on Quantum Logic}. \emph{arXiv:2610.34266}  (2026).
\bibitem{indirect} S. Vedl, V. Srivastava, R. Ghosh, A. Gilchrist. \href{https://arxiv.org/abs/2609.18409}{Benchmarking indirect quantum control schemes via higher-order quantum operations}. \emph{arXiv:2609.18409}  (2026).
\bibitem{qsp} G. H. Low, I. L. Chuang. \href{https://doi.org/10.1103/PhysRevLett.118.010501}{Optimal Hamiltonian Simulation by Quantum Signal Processing}. \emph{Physical Review Letters} 118 010501 (2017).
\bibitem{qubitization} G. H. Low, I. L. Chuang. \href{https://doi.org/10.22331/q-2019-07-12-163}{Hamiltonian Simulation by Qubitization}. \emph{Quantum} 3 163 (2019).
\bibitem{qsvt} A. Gilyén, Y. Su, G. H. Low, N. Wiebe. \href{https://doi.org/10.1145/3313276.3316366}{Quantum singular value transformation and beyond: exponential improvements for quantum matrix arithmetics}. \emph{Proceedings of the 51st Annual ACM SIGACT Symposium on Theory of Computing} 193-204 (2019).
\bibitem{grand} J. M. Martyn, Z. M. Rossi, A. K. Tan, I. L. Chuang. \href{https://doi.org/10.1103/PRXQuantum.2.040203}{Grand Unification of Quantum Algorithms}. \emph{PRX Quantum} 2 040203 (2021).
\bibitem{qn} X. Lu, Y. Liu, H. Lin. \href{https://doi.org/10.22331/q-2026-03-27-2048}{Quantum Signal Processing and Quantum Singular Value Transformation on U(N)}. \emph{Quantum} 10 2048 (2026).
\bibitem{mqsp} G. H. Low. \href{https://arxiv.org/abs/2610.01125}{Multivariate Quantum Signal Processing}. \emph{arXiv:2610.01125}  (2026).
\bibitem{linearqsp} M. Arsenault, H. Kristjánsson. \href{https://arxiv.org/abs/2608.14387}{Linearised quantum signal processing}. \emph{arXiv:2608.14387}  (2026).
\bibitem{arbmatrix} X. Gutiérrez, L. Laneve, M. Sanz. \href{https://arxiv.org/abs/2604.19688}{Quantum Eigenvalue Transformations for Arbitrary Matrices}. \emph{arXiv:2604.19688}  (2026).
\bibitem{recentqsp} T. Gurfinkel, K. De Silva, A. Mahasinghe, et al. \href{https://arxiv.org/abs/2609.09977}{From Block-encoding to Generalized Quantum Signal Processing: Principles, Algorithms and Applications}. \emph{arXiv:2609.09977}  (2026).
\bibitem{trotter26} G. Rendón. \href{https://doi.org/10.22331/q-2026-11-06-2261}{All you need is Trotter}. \emph{Quantum} 10 2261 (2026).
\bibitem{schmidt} E. Schmidt. \href{https://doi.org/10.1007/BF01449770}{Zur Theorie der linearen und nichtlinearen Integralgleichungen}. \emph{Mathematische Annalen} 63 433-476 (1907).
\bibitem{cleve} R. Cleve. \href{https://cleve.iqc.uwaterloo.ca/resources/Qic890LectureNotes2019Apr22\%28V22\%29.pdf}{Introduction to Quantum Information Processing}. \emph{University of Waterloo, QIC 890 lecture notes}  (2019).
\bibitem{tensorrank} V. de Silva, L. Lim. \href{https://doi.org/10.1137/06066518X}{Tensor Rank and the Ill-Posedness of the Best Low-Rank Approximation Problem}. \emph{SIAM Journal on Matrix Analysis and Applications} 30 1084-1127 (2008).
\bibitem{svm} C. Cortes, V. Vapnik. \href{https://doi.org/10.1007/BF00994018}{Support-vector networks}. \emph{Machine Learning} 20 273-297 (1995).
\bibitem{chen} Y. Chen, W. Huang, L. M. Nguyen, T.-W. Weng. \href{https://proceedings.neurips.cc/paper/2021/hash/c559da2ba967eb820766939a658022c8-Abstract.html}{On the Equivalence between Neural Network and Support Vector Machine}. \emph{NeurIPS 34}  (2021).
\bibitem{ntk} A. Jacot, F. Gabriel, C. Hongler. \href{https://arxiv.org/abs/1806.07572}{Neural Tangent Kernel: Convergence and Generalization in Neural Networks}. \emph{NeurIPS 31}  (2018).
\bibitem{gunasekar} S. Gunasekar, J. D. Lee, D. Soudry, N. Srebro. \href{https://proceedings.neurips.cc/paper/2018/hash/0e98aeeb54acf612b9eb4e48a269814c-Abstract.html}{Implicit Bias of Gradient Descent on Linear Convolutional Networks}. \emph{NeurIPS 31}  (2018).
\bibitem{cybenko} G. Cybenko. \href{https://doi.org/10.1007/BF02551274}{Approximation by superpositions of a sigmoidal function}. \emph{Mathematics of Control, Signals, and Systems} 2 303-314 (1989).
\bibitem{noprogram} M. A. Nielsen, I. L. Chuang. \href{https://doi.org/10.1103/PhysRevLett.79.321}{Programmable Quantum Gate Arrays}. \emph{Physical Review Letters} 79 321-324 (1997).
\bibitem{tatsu} TatSu project. \href{https://tatsu.readthedocs.io/en/stable/}{TatSu documentation}. \emph{Stable documentation, accessed November 8}  (2026).
\bibitem{zen} T. Peters. \href{https://peps.python.org/pep-0020/}{PEP 20 -- The Zen of Python}. \emph{Python Enhancement Proposals}  (2004).
\bibitem{cupy} CuPy project. \href{https://docs.cupy.dev/en/stable/user_guide/kernel.html}{CuPy user-defined kernels}. \emph{Documentation, accessed November 8}  (2026).
\bibitem{lean} Lean project. \href{https://lean-lang.org/doc/reference/latest/}{The Lean Language Reference}. \emph{Documentation, accessed November 8}  (2026).
\bibitem{braket} Amazon Web Services. \href{https://docs.aws.amazon.com/braket/latest/developerguide/what-is-braket.html}{Amazon Braket Developer Guide}. \emph{Documentation, accessed November 8}  (2026).
\bibitem{cirq} Google Quantum AI. \href{https://quantumai.google/cirq}{Cirq}. \emph{Documentation, accessed November 8}  (2026).
\bibitem{azure} Microsoft. \href{https://learn.microsoft.com/en-us/azure/quantum/provider-rigetti}{Rigetti provider}. \emph{Azure Quantum documentation, accessed November 8}  (2026).
\bibitem{dwave} D-Wave. \href{https://docs.dwavequantum.com/en/latest/ocean/}{Ocean software}. \emph{Documentation, accessed November 8}  (2026).
\bibitem{eh} M. Ekerå, J. Håstad. \href{https://arxiv.org/abs/1702.00249}{Quantum algorithms for computing short discrete logarithms and factoring RSA integers}. \emph{arXiv:1702.00249}  (2017).
\bibitem{shortpost} M. Ekerå. \href{https://doi.org/10.1007/s10623-020-00783-2}{On post-processing in the quantum algorithm for computing short discrete logarithms}. \emph{Designs, Codes and Cryptography} 88 2313-2335 (2020).
\bibitem{shortsuccess} M. Ekerå. \href{https://arxiv.org/abs/2309.01754}{On the success probability of the quantum algorithm for the short DLP}. \emph{arXiv:2309.01754}  (2023).
\bibitem{orderpost} M. Ekerå. \href{https://doi.org/10.1145/3655026}{On the Success Probability of Quantum Order Finding}. \emph{ACM Transactions on Quantum Computing} 5 1-40 (2024).
\bibitem{quaspy} M. Ekerå. \href{https://github.com/ekera/quaspy}{Quaspy: quantum algorithm simulation in Python}. \emph{Reference implementation, version 0.9.4}  (2026).
\bibitem{gidney} C. Gidney. \href{https://arxiv.org/abs/2505.15917}{How to factor 2048 bit RSA integers with less than a million noisy qubits}. \emph{arXiv:2505.15917}  (2025).
\bibitem{pinnacle} P. Webster, L. Berent, O. Chandra, et al. \href{https://arxiv.org/abs/2602.11457}{The Pinnacle Architecture: Reducing the cost of breaking RSA-2048 to 100 000 physical qubits using quantum LDPC codes}. \emph{arXiv:2602.11457}  (2026).
\bibitem{oratomic} M. Cain, Q. Xu, R. King, et al. \href{https://arxiv.org/abs/2603.28627}{Shor's algorithm is possible with as few as 10,000 reconfigurable atomic qubits}. \emph{arXiv:2603.28627}  (2026).
\bibitem{catshor} A. Pappalardo, D. Goudarzi, V. Hertz, et al. \href{https://arxiv.org/abs/2610.37839}{Shor's algorithm compiled on cat qubits: 22 593 qubits for RSA-2048 in 3.14 days}. \emph{arXiv:2610.37839}  (2026).
\bibitem{torch} PyTorch project. \href{https://docs.pytorch.org/docs/stable/generated/torch.nn.Sequential.html}{Sequential}. \emph{PyTorch documentation}  (2026).
\bibitem{shende} V. V. Shende, S. S. Bullock and I. L. Markov. \href{https://doi.org/10.1109/TCAD.2005.855930}{Synthesis of quantum-logic circuits}. \emph{IEEE Transactions on Computer-Aided Design of Integrated Circuits and Systems} 25, 1000--1010 (2006).
\end{thebibliography}
